% Source: https://artofproblemsolving.com/community/c3h1405643_counting_and_probability
% Stable identity: aops-topic-1405643-q01
\textbf{(a)} A deal is an ordered pair (Amelia's hand, John's hand) of disjoint $7$-card sets, and all such pairs are equally likely. Swapping the two hands is a bijection from the set of deals onto itself, and it maps the deals in which Amelia's sum is larger exactly onto the deals in which John's sum is larger. So these two events have the same probability $p$. The three events ``Amelia larger'', ``John larger'', ``equal'' are disjoint and cover all deals, so $2p+q=1$ and $p=\dfrac{1-q}{2}$.

\textbf{(b)} It suffices to show $q>0$, that is, that equal sums can happen. For example, Amelia may get $\{1,2,3,4,5,6,49\}$ and John $\{7,8,9,10,11,12,13\}$; both sums are $70$. Since this deal has positive probability, $q>0$ and so $p=\dfrac{1-q}{2}<\dfrac{1}{2}$.

So the poster's guess $\dfrac{1}{2}$ is not correct; the true value is slightly less than $\dfrac{1}{2}$ (the difference is half the probability of a tie).
